\input{../tex-inputs/latex-leseansicht-vorspann}

\section[Gustav Schwarzkopf an Arthur Schnitzler, 19. 8. 1926]{L04253 Gustav Schwarzkopf an Arthur Schnitzler, 19. 8. 1926}
\nopagebreak\mylabel{L04253v}
\rehead{ }\normalsize\beginnumbering\briefempfaengerindex{Schnitzler, Arthur@\textsc{Schnitzler, Arthur}!zzzSchwarzkopf, Gustav@\emph{von Gustav Schwarzkopf}!1926-08-191@{19. 8. 1926}|(be}
\toendnotes[C]{\smallbreak\pagebreak[2]}
\correspDesc{Versand  durch Gustav Schwarzkopf am 19. 8. 1926 in Wien
\newline{}Erhalt  durch Arthur Schnitzler im Zeitraum [20. 8. 1926 – 24. 8. 1926?] in Adelboden}\toendnotes[C]{\smallbreak}
\Standort{CUL, Schnitzler, B 96a.}
\physDesc{Postkarte, 1089 Zeichen
\newline{}Handschrift: schwarze Tinte, deutsche Kurrent
\newline{}Versand: 1) Stempel: »\nobreak{}\oindex{I., Innere Stadt@\textbf{I., Innere Stadt}, \emph{Verwaltungsgebiet}|pwk}1/1 Wien 8, 9. VIII. 26, \textcolor{gray}{1}8\nobreak{}«.   2) Stempel: »\nobreak{}Besuchet die Wiener Messe 3.–12. Sept. 1926\nobreak{}«. }\toendnotes[C]{\smallbreak}\pstart{}{\pb}I. Tiefer Graben 17\oindex{Wien@\textbf{Wien}!I., Innere Stadt@\textbf{I., Innere Stadt}!Tiefer Graben 17@\textbf{Tiefer Graben 17}, \emph{Wohngebäude}|pw}\pend{}\pstart{}II. 6\pend{}{\bigskip}\pstart{}Herrn\pend{}\pstart{}\textsc{D\textsuperscript{r} Arthur Schnitzler}\pend{}\pstart{}\textsc{Adelboden\oindex{Adelboden@\textbf{Adelboden}|pw}}\pend{}\pstart{}\textsc{(Bern\oindex{Bern@\textbf{Bern}, \emph{Hauptstadt}|pw})}\pend{}\pstart{}\textsc{(Schweiz\oindex{Schweiz@\textbf{Schweiz}|pw})}\pend{}\pstart{}\textsc{Hotel Melada\oindex{Nevada Palace@\textbf{Nevada Palace}, \emph{Hotel}|pw}}\pend{}{\bigskip}\vspace{1em}
\pstart
           \raggedleft{}{\pb}9. VIII 26\pend
           \vspace{0.5em}
\pstart
           Lieber Arthur,{ }ſie habe mich aufgefordert über \textsc{Emil\pwindex{Schwarzkopf, Emil 17.\,9.\,1851 Wien – 28.\,1.\,1928 ebd.@\textsc{Schwarzkopf, Emil} (17.\,9.\,1851 Wien – 28.\,1.\,1928 ebd.), \emph{Übersetzer, Komponist}|pw}} Nachricht zu geben. Bis geſtern war nichts zu melden. Heute hat mich der A. Dr
                  \textsc{Jesc{[}h{]}ke\pwindex{Jeschke, Karl @\textsc{Jeschke, Karl}, \emph{Hautarzt, Mediziner}|pw}} rufen laſſen und mir mitgeteilt: der Ausſchlag{ }ſei{ }ſo gut wie geheilt, aber man
               habe die Abſicht den Kranken\pwindex{Schwarzkopf, Emil 17.\,9.\,1851 Wien – 28.\,1.\,1928 ebd.@\textsc{Schwarzkopf, Emil} (17.\,9.\,1851 Wien – 28.\,1.\,1928 ebd.), \emph{Übersetzer, Komponist}|pwv}
               doch noch einige Wochen im Hauſe zu behalten. Etwa bis Ende Auguſt. Da{\geminationn} müſſen, um Reinigungsarbeiten vorzunehmen, die Kranken
               dieſes Zi{\geminationm}ers auf die andern Krankenzi{\geminationm}er verteilt werden, es iſt alſo nicht mehr viel
               überflüſſiger Raum vorhanden. Er hat mich gefragt, ob ich einverſtanden bin, we{\geminationn} von Seite des Spitals\orgindex{Wiedner Spital@Wiedner Spital|pwv} wegen Aufnahme in’s \label{K_L04253-1v}\edtext{Verſorgungshaus}{\lemma{\textnormal{\emph{Versorgungshaus}}}\Cendnote{\textnormal{Armenhaus der öffentlichen Hand}}}\label{K_L04253-1}
               intervenirt wird. Ich habe für alle Fälle zugeſti{\geminationm}t und
               werde mich jetzt noch um eine eventuelle andere Unterkunft umſehen: – Sonſt nichts
               Neues. – Hoffentlich haben Sie beſſeres Wetter als wie hier. Seit 8 Wochen{ }ſteht Max\pwindex{Schwarzkopf, Max 12.\,6.\,1857 Wien – 14.\,4.\,1928 ebd.@\textsc{Schwarzkopf, Max} (12.\,6.\,1857 Wien – 14.\,4.\,1928 ebd.), \emph{Rechtsanwalt}|pw}’ens Reisetaſche gepackt, aber{ }ſelbſt für
               die {\pb}Reiſe nach Baden\oindex{Baden bei Wien@\textbf{Baden bei Wien}, \emph{Hauptstadt}|pw} war das Wetter nicht gerad ganz
               verlockend.\pend
           
\pstart
           Viele Grüße an Frau \textsc{Olga\pwindex{Schnitzler, Olga 17.\,1.\,1882 Wien – 13.\,1.\,1970 Lugano@\textsc{Schnitzler, Olga} (17.\,1.\,1882 Wien – 13.\,1.\,1970 Lugano), \emph{Schauspielerin, Sängerin}|pw}}, \textsc{Heini\pwindex{Schnitzler, Heinrich 9.\,8.\,1902 Hinterbrühl – 12.\,7.\,1982 Wien@\textsc{Schnitzler, Heinrich} (9.\,8.\,1902 Hinterbrühl – 12.\,7.\,1982 Wien), \emph{Regisseur, Schauspieler}|pw}} (iſt nicht \label{K_L04253-2v}\edtext{heute{ }ſein Geburtstag}{\lemma{\textnormal{\emph{heute sein Geburtstag}}}\Cendnote{\textnormal{Ja, Heinrich Schnitzler\pwindex{Schnitzler, Heinrich 9.\,8.\,1902 Hinterbrühl – 12.\,7.\,1982 Wien@\textsc{Schnitzler, Heinrich} (9.\,8.\,1902 Hinterbrühl – 12.\,7.\,1982 Wien), \emph{Regisseur, Schauspieler}|pwk} wurde am 9. 8. 1926 24 Jahre alt.}}}\label{K_L04253-2}?){[},{]}{ }\textsc{Lili\pwindex{Cappellini, Lili 13.\,9.\,1909 Wien – 26.\,7.\,1928 Venedig@\textsc{Cappellini, Lili} (13.\,9.\,1909 Wien – 26.\,7.\,1928 Venedig)|pw}}\pend
           
\pstart
           Herzlichſt{\\[\baselineskip]}Ihr{\\[\baselineskip]}\spacefill\mbox{Gustav.}\pend
           \leftskip=0em{}\selectlanguage{ngerman}\endnumbering\briefempfaengerindex{Schnitzler, Arthur@\textsc{Schnitzler, Arthur}!zzzSchwarzkopf, Gustav@\emph{von Gustav Schwarzkopf}!1926-08-191@{19. 8. 1926}|)be}\mylabel{L04253h}
\begin{anhang}
\end{anhang}\newcommand{\dateiname}{L04253}\newcommand{\titel}{Gustav Schwarzkopf an Arthur Schnitzler, 19. 8. 1926}\newcommand{\editorInnen}{Herausgegeben von Jahnke, SelmaMüller, Martin Anton}\input{../tex-inputs/latex-leseansicht-abspann}
